\documentclass[10pt]{article}
% Shared look for overview.tex and project-details.tex. Build with: tectonic -X compile <file>.tex
\usepackage[a4paper,margin=1.7cm]{geometry}
\usepackage{fontspec}
\setmainfont{texgyrepagella}[Extension=.otf,UprightFont=*-regular,BoldFont=*-bold,ItalicFont=*-italic,BoldItalicFont=*-bolditalic]
\setsansfont{texgyreheros}[Extension=.otf,UprightFont=*-regular,BoldFont=*-bold,ItalicFont=*-italic,BoldItalicFont=*-bolditalic]
\usepackage{amsmath,microtype,xcolor,tikz,booktabs,tabularx,enumitem,titlesec,array}
\usepackage[hidelinks]{hyperref}
\usetikzlibrary{positioning,arrows.meta,fit}

\definecolor{ink}{HTML}{1F2933}
\definecolor{accent}{HTML}{1F6F8B}
\definecolor{soft}{HTML}{E8F1F4}
\definecolor{muted}{HTML}{5C6B73}
\color{ink}

\titleformat{\section}{\sffamily\bfseries\large\color{accent}}{}{0pt}{}
\titlespacing*{\section}{0pt}{1.1ex}{0.5ex}
\titleformat{\subsection}{\sffamily\bfseries\normalsize}{}{0pt}{}
\titlespacing*{\subsection}{0pt}{1ex}{0.3ex}
\setlist{itemsep=1pt,topsep=2pt,parsep=0pt,leftmargin=1.2em}
\setlength{\parindent}{0pt}
\setlength{\parskip}{0.5ex}
\pagestyle{empty}

\newcommand{\code}[1]{\texttt{\small #1}}
% \gh[branch]{path}{text}: a link to that place in the group repo on GitHub (default branch main)
\newcommand{\repo}{https://github.com/angelabuxo/IMDb-movie-success}
\newcommand{\gh}[3][main]{\href{\repo/tree/#1/#2}{\color{accent}#3}}
\newcommand{\where}[1]{{\small\color{muted}#1}}
\newcolumntype{L}{>{\raggedright\arraybackslash}X}

\newcommand{\doctitle}[2]{%
  {\sffamily\bfseries\LARGE\color{accent}#1\par}\vspace{2pt}
  {\small\color{muted}#2\par}\vspace{4pt}}


\begin{document}
\doctitle{What makes a movie profitable? Project overview}{Log-Log Legends: Àngela Buxó, Gyula Kürthy, Niklas Johansen · DTU 02805, autumn 2026 · state on 14 November · paper due 27 November, presentation 2 December · more detail in \gh[g/project-overview]{docs/overview/project-details.pdf}{project-details.pdf} · \gh{}{GitHub repo}}

\textbf{Question.} What makes a movie highly profitable, and does making money agree with being liked?

\textbf{Profit} is $\ln(\text{worldwide gross}/\text{budget})$, ``the multiple''. It leaves out marketing and the cinemas' share, so
it's a return on the reported budget, not real profit. \textbf{Sample:} 2,446 clean films with 100,000+ IMDb votes, or
6,851 with 10,000+.

\section{Three strands, one model}

\begin{center}
\begin{tikzpicture}[
  box/.style={draw=accent, fill=soft, rounded corners=3pt, align=center, font=\small, inner sep=5pt, text width=4.6cm, minimum height=1.7cm},
  head/.style={font=\sffamily\bfseries\small},
  arr/.style={-{Stealth[length=5pt]}, accent, thick},
  node distance=6mm and 5mm]
  \node[box, text width=15.6cm, minimum height=0pt] (data) {\textbf{Data}: IMDb · Wikipedia · Wikidata · MovieSum screenplays · MovieLens Tag Genome reviews};
  \node[box, below=of data] (text) {{\sffamily\bfseries What it says}\\ word comparison (Scattertext)\\ emotional arcs\\ script features};
  \node[box, left=of text] (net) {{\sffamily\bfseries Who made it}\\ movie graph\\ actor network\\ shared-people returns};
  \node[box, right=of text] (rec) {{\sffamily\bfseries How it was received}\\ critics vs audience vs money\\ review tone and words};
  \node[box, below=of text, text width=15.6cm, minimum height=0pt] (model) {\textbf{What predicts profit}: every block side by side, trained on films up to 2010, tested on 2011 to 2017};
  \foreach \n in {net,text,rec} {\draw[arr] (data.south -| \n) -- (\n.north); \draw[arr] (\n.south) -- (model.north -| \n);}
\end{tikzpicture}
\end{center}

\section{What we've found}

\begin{tabularx}{\linewidth}{@{}>{\bfseries}p{3.2cm}L@{}}
\toprule
Who made it & The best predictor known before release. Studio and people track records lift out-of-sample $R^2$ from $-0.01$ to $0.20$. Films that share people have similar returns. \\
\midrule
What it says & Almost nothing once budget, genre and era are fixed. The word difference vanishes without franchise films, and the arcs show no clear link (one near miss). \\
\midrule
How it was received & Money and liking agree only moderately ($\rho$ 0.27 to 0.33). Reviews add to the model (\mbox{$R^2$ 0.24 to 0.31}), but they come after release. \\
\bottomrule
\end{tabularx}

\medskip
\textbf{Possible one story:} who makes a film predicts its return, what's in the script doesn't, and audiences
only partly reward what earns.

\section{Rules every strand follows}
\begin{itemize}
  \item Network measures use only earlier films.
  \item Every result gets a check: a shuffle null, a second method, or hand labels.
  \item Budget, era and genre are controlled, and franchise films get a separate check.
  \item No script text, review text or IMDb fields go into git.
\end{itemize}

\section{Open}
\begin{itemize}
  \item Agree on the one story for the paper.
  \item Fill the human validation sheets (Scattertext, arcs, reviews).
  \item Merge Àngela's \gh[reviews-exploration]{analysis/reviews}{reviews branch} and Niklas's \gh[initial_actor_network]{scripts/network/actor_roi.py}{actor-network branch}.
  \item Arcs Step 11 (script tone vs review tone) needs the group to decide whether whole-review VADER counts as the same instrument.
\end{itemize}
\end{document}
